%- 导言区开始 - by Granvallen
\documentclass[a4paper,10.5pt,oweside]{article} % 注意单页时是不分奇偶页的 
% twoside 双页 oweside 单页

    %- Font -
    \usepackage[slantfont,boldfont,AutoFakeBold=true]{xeCJK} % XeCJK中文排版支持 
        % AutoFakeBold=false % 能解决pdf复制出来乱码的问题 但问题是这样的话要自己设置粗体才能用粗体(原来默认是提供一个伪粗体显示) 可以通过\bfseries命令使之变成粗体
        % 所以要么使用黑体来代替 要么单独找 粗宋体字体. 另外的解决办法是用 Adobe Song Std 代替 宋体
        \setmainfont{Times New Roman} % 设置英文字体
        % \setCJKmonofont{Adobe Song Std} % 设置缺省中文字体
        \setCJKmainfont{SimSun} % 设置中文字体 宋体  YouYuan  NSimSun
        \newCJKfontfamily\YouYuan{YouYuan} % 设置其他中文字体命令 幼圆
        \newCJKfontfamily\SimSun{SimSun}
        \xeCJKEditPunctStyle{quanjiao}{bound-margin-ratio=0.25} % 符号缩进
        %\parindent = 2em   % 段首缩进  % 后面使用 indentfirst 宏包实现缩进比较好
    \usepackage{fontspec} % 字体命令
       \newfontfamily\monaco{Monaco} % 这里英语的字体命令是有用的, 中文的就不行 可能是被XeCJK覆盖了

    %- Font CMD -
    % 常用命令
    % \bfseries{} -> 加粗
    % \slshape{} -> 斜体
    % \boldsymbol{} -> 粗体字母 表示向量与矩阵
    % 字体大小命令 使用 {\chuhao XXX}
    \newcommand{\chuhao}{\fontsize{42pt}{\baselineskip}\selectfont}
    \newcommand{\xiaochuhao}{\fontsize{36pt}{\baselineskip}\selectfont}
    \newcommand{\yihao}{\fontsize{28pt}{\baselineskip}\selectfont}
    \newcommand{\erhao}{\fontsize{21pt}{\baselineskip}\selectfont}
    \newcommand{\xiaoerhao}{\fontsize{18pt}{\baselineskip}\selectfont}
    \newcommand{\sanhao}{\fontsize{15.75pt}{\baselineskip}\selectfont}
    \newcommand{\sihao}{\fontsize{14pt}{\baselineskip}\selectfont}
    \newcommand{\xiaosihao}{\fontsize{12pt}{\baselineskip}\selectfont}
    \newcommand{\wuhao}{\fontsize{10.5pt}{\baselineskip}\selectfont}
    \newcommand{\xiaowuhao}{\fontsize{9pt}{\baselineskip}\selectfont}
    \newcommand{\liuhao}{\fontsize{7.875pt}{\baselineskip}\selectfont}
    \newcommand{\qihao}{\fontsize{5.25pt}{\baselineskip}\selectfont}

    %- Math -
    \usepackage{amsmath, amsthm, amssymb} % 数学支持
        % 定理环境风格
        \theoremstyle{definition}   
        %\theoremstyle{plain}
        %\theoremstyle{remark}
        % 定理 定义 性质 
        \newtheorem{definition}{定义}[section] % 定义环境
        \newtheorem{theorem}{定理}[section] % 定理环境
        \newtheorem{property}{性质}[section] % 性质环境
        % 证明环境
        % \newenvironment{pf}{{\noindent\bf 证明:}\quad}{\hfill $\qed$\par}  
        \newenvironment{pf}[1]{\par \textbf{证明}\quad #1}{\hfill$\qed$ \par}
        \renewcommand*{\qedsymbol}{[证毕]}
        % 公式编号
        \numberwithin{equation}{section} % 公式编号随章节走
    % 其他数学字体支持
    \usepackage{mathrsfs} % 字母花体
    \usepackage{dsfont} % 漂亮的R
    % 算法伪代码格式支持
    \usepackage{algorithm}
    \usepackage{algorithmic}
    % 矩阵格式调整
    \makeatletter
    \def\env@matrix{\hskip -\arraycolsep
    \let\@ifnextchar\new@ifnextchar
    \linespread{1}\selectfont
    \renewcommand{\arraystretch}{1.2}%
    \array{*\c@MaxMatrixCols c}}
    \makeatother

    %- Page -
    \usepackage{geometry}
        \geometry{
            paper=a4paper,
            left=2.7cm,   % 页边距设定
            right=2.3cm,
            top=2.2cm, 
            bottom=2cm, 
            headsep=0.7cm, % 页眉到正文距离
            foot=0.7cm, % 页脚到正文距离
        }
    \usepackage{titlesec} % 设置标题显示格式 一级标题 二级标题
        % \titleformat{\part}{\centering\sanhao\bfseries}{\thepart}{1em}{}% 设置part格式 居中 字体3号 加粗 
        \titleformat{\section}{\xiaosihao\bfseries}{\thesection}{1em}{}
        \titleformat{\subsection}{\wuhao\bfseries}{\thesubsection}{1em}{}
        \titleformat{\subsubsection}{\wuhao\bfseries}{\thesubsubsection}{1em}{}
        \titleformat{\paragraph}{\wuhao\bfseries}{\theparagraph}{1em}{}
    \usepackage{titletoc} % 定制目录显示格式
        % \titlecontents{part}[0em] % 标签距离页面左边距
        %     {\xiaosihao\bfseries} % 字体
        %     {\contentslabel{1em}}  % 标签距离目录内容距离
        %     {}
        %     {\titlerule*{$\dots$}\contentspage\hspace*{0em}} % 指引线与页码
        % \titlecontents{section}[2em] % 标签距离页面左边距
        %     {\xiaosihao} % 字体
        %     {\contentslabel{1em}}  % 标签距离目录内容距离
        %     {}
        %     {\titlerule*{$\cdots$}\contentspage\hspace*{1em}} % 指引线与页码
        % \titlecontents{subsection}[3em] % 标签距离页面左边距
        %     {\xiaosihao} % 字体
        %     {\contentslabel{2em}}  % 标签距离目录内容距离
        %     {}
        %     {\titlerule*{$\cdots$}\contentspage\hspace*{1em}} % 指引线与页码
        % \titlecontents{subsubsection}[4em] % 标签距离页面左边距
        %     {\xiaosihao} % 字体
        %     {\contentslabel{3em}}  % 标签距离目录内容距离
        %     {}
        %     {\titlerule*{$\cdots$}\contentspage\hspace*{1em}} % 指引线与页码
    \usepackage{fancyhdr} % 页眉页脚设定
        % \pagestyle{fancy}
        % \fancyhf{}
        % 页眉页脚分割线 
        % \renewcommand{\headrulewidth}{0.4pt}
        % \renewcommand{\footrulewidth}{0.4pt}
        % 页眉页脚设定 有两种方法
        % 1
        % \fancyhead[OC]{\wuhao XX大学学士学位论文} % OC 表示设定奇数页的中间 默认是偶数页
        % \fancyhead[C]{\wuhao 基于XX}
        % \fancyfoot[C]{\wuhao\thepage} 
        % 2  这种方法好像无法设定奇偶页区分显示
        % \lhead{} % 页眉左侧内容
        % \chead{\wuhao XX大学学士学位论文}
        % \rhead{} 
        % \lfoot{}
        % \cfoot{}
    
    %- Paragraph -
    \usepackage{indentfirst} % 设置首行缩进
        \setlength{\parindent}{2em}
    \linespread{1.2} % 行距
    % \setlength{\columnsep}{1em} %列间距  好像不管用
    % \setlength{\topmargin}{10mm} % 页眉到页边的距离
    % \setlength{\parskip}{10pt}  % 设置部分段落间隔
    \def\changemargin#1#2{\list{}{\rightmargin#2\leftmargin#1}\item[]}
    \let\endchangemargin=\endlist % 增加可更改段落页边距 

    %- TIKZ -
    \usepackage{tikz,pgfplots} % 画图支持    如用 matlab2tikz 的输出需要 pgfplots
        \usetikzlibrary{graphs,arrows,shapes,chains,quotes,plotmarks}
    \pgfplotsset{compat=1.16}
    \newenvironment{flowchart}{ % 一个画流程图环境
        \begin{tikzpicture}[
        %inner sep=2mm, % 边框到形状内字符距离
        > = stealth, % 箭头样式
        -> /.style={thick ,black},
        -- /.style={thick ,black},
        line width=1pt,
        every edge quotes/.style={font=\sihao,auto}, % 标签样式
        every node/.style={font=\sihao}, % 框内字体
        point/.style={inner sep=0pt,line width=0pt,fill}, % point样式
        terminal/.style={circle,draw,thin,inner sep=0pt}, % 结点样式
        rect/.style={rectangle,rounded corners=1mm,draw,thick,inner sep=5pt,minimum size=10mm,align=center}, %rect样式
        operation/.style={circle,draw,thick,inner sep=-1pt}, % operation样式
        hv path/.style={to path={-| (\tikztotarget)}}, % 直角折线
        vh path/.style={to path={|- (\tikztotarget)}},
        ]}{\end{tikzpicture}}
    
    %- Table -
    \usepackage{multirow}
    \usepackage{booktabs} % 三线表支持
    \usepackage{longtable} % 跨页表格支持 
    \renewcommand\arraystretch{1.5} % 表格行距
    \renewcommand\tabcolsep{3mm} % 表格列距
    \renewcommand\tablename{表} % 改表的标号

    %- Picture -
    \usepackage{graphicx} % 插入图片
    \usepackage{subfigure} % 插入并列图
    \usepackage{pdfpages} % 插入pdf支持
    \usepackage{caption} % 去掉表/图的标号需要
        \captionsetup[table]{labelsep=quad} % 去掉 表几 后面的冒号
        \captionsetup[figure]{labelsep=quad} % 去掉 图几 后面的冒号
        \setlength{\abovecaptionskip}{10pt} % 设定图与caption距离
        \setlength{\belowcaptionskip}{-10pt}
        \DeclareCaptionFormat{captionformat}{\wuhao#1#2#3} % 改图表标签字号
        \captionsetup{format=captionformat}
    \usepackage{wrapfig} % 字体环绕图片支持
    \usepackage{float} % 插入图片可使用 'H' 取消浮动体机制
    \usepackage{xcolor} % 颜色支持 用于代码高亮 和 tikz绘图
    \renewcommand\figurename{图} % 改图的标号
    \usepackage{epstopdf} % 提供eps图片插入支持    matlab2tikz 的实现并不完美 还是插入eps比较好用

    %- Code -
    \usepackage{listings} % 代码环境
        \lstset{
            showstringspaces=false,
            columns=fixed,
            numbers=left,
            frame=none,
            backgroundcolor=\color[RGB]{245,245,220},
            keywordstyle=\color[RGB]{0,0,205},
            %numberstyle=\footnotesize\color{darkgray},
            commentstyle=\it\color[RGB]{0,96,96},
            stringstyle=\rmfamily\slshape\monaco\color[RGB]{47,79,79},
            numberstyle=\tiny\monaco,
            basicstyle=\tiny\monaco,
            escapeinside=``,
            language=c++,
        }
    
    
    %- Bibliography -
    % 参考文献角标设定
    % 角标设定为: 数字编号、排序且压缩、上标、外侧方括    压缩是说当引用多与三个时使用 '-' 连接
    \usepackage[numbers, sort&compress, super, square]{natbib} 
    \renewcommand{\refname}{\bfseries\xiaosihao{参考文献}} % article使用这个
    %\renewcommand{\bibname}{参考文献} % book使用这个
    \renewcommand{\bibfont}{\xiaowuhao} % 参考文献项的字体
    % \renewcommand{\bibnumfmt}[1]{\YouYuan{[#1]}} % 修改文献编号样式
        \setlength{\bibsep}{0.25ex} % 文献条目之间的距离    单位 倍行距
        % \bibliographystyle{unsrt} % 设置参考文献排序的类型 按引用顺序 用gbt7714就不需要这个了
    \usepackage{gbt7714} % 国标文献格式支持

    % 附录
    \usepackage{appendix}

    % 引用
    % \usepackage{showkeys} % 写的时候开启 提示引用label 方便差错

    %- Link - 放最后hyperref易与其他宏包冲突
    \usepackage[hidelinks]{hyperref} % 超链接 与 书签 支持    提供\phantomsection

    %- 默认封面信息 -
    % \title{高斯过程回归方法及其应用}
    % \author{作者: XX}
    % \date{2018.11.1}% 默认使用今天的日期  怎么写怎么打印

    % 自定义命令
    \newcommand{\Title}[1]{\vspace{1ex} {\sanhao \bfseries #1}}
    \newcommand{\AuthorID}[2]{\vspace{3.5ex} {\xiaosihao \bfseries #2,\: #1}}
    \newcommand{\Date}[1]{\vspace{2ex} {\wuhao #1} \vspace{2ex}}
    \newcommand{\Abstractcn}[2]{
        \begin{changemargin}{1cm}{1cm}
            {\xiaowuhao{\bfseries{摘\quad 要:\ }} #1 } \\ \vspace{1ex}
            {\xiaowuhao {\bfseries 关键词:} #2 }
        \end{changemargin}}
    \newcommand{\Abstracten}[2]{
        \begin{changemargin}{1cm}{1cm}
            {\xiaowuhao{\bfseries{Abstract:}} #1 } \\ \vspace{1ex}
            {\xiaowuhao {\bfseries Key words:} #2 }
        \end{changemargin}}



%- 导言区结束 -

\begin{document}
    %- 封面与目录 -

    % - 前置 -
    \includepdf[pages=1]{./封面与尾页/封面.pdf} % pdf格式的封面 可以用pages指定页数


    %- 正文开始 - -------------------------------------------------------
    \setcounter{page}{1}
    %- 题头信息 -
    \leavevmode % 使第一行的\vspace生效
    \begin{center}
        \vspace{1ex}

        \Title{题目三号粗体居中，段前空1行，段后空0.5行}

        \AuthorID{姓名（小四粗体居中，段前后各0.5行）}{学号}

        \Date{日期（五号居中，段前0.5行，段后1行） }

        % 中文摘要
        \Abstractcn{这是文献报告写作要求。排版格式要求包括用纸，页面设置、字体、行距等。基本要求中强调独立撰写，不可简单剪贴现成资料。对文献的数量和质量作出规定，提供文献清单实例，指出一些通病，并给出评分依据。\textcolor{red}{注意中英文摘要一致。}}{}

        \vspace{3ex}

        \Title{English Title Uses the Same Format as the Chinese Title}

        % 英文摘要
        \Abstracten{An English abstract must be included, which should be 400~500 words in length. This will be a part of the basis for the assessment of your technical English proficiency. You must write it in your own words, rather than just copy from a reference. Mind you that anyone who simply copies from materials written by others will not pass the examination. \textcolor{red}{（200 字左右，与以上中文摘要一致）}}{}

    \end{center}
    \vspace{2ex}



    % %- 章节 -
    % % !TEX root = ../thesis.tex
\newpage

\section{前言}
    XXX































    % !TEX root = ../thesis.tex

\section{引言}
    文献报告格式要求：用A4纸，通栏排版。页边距：上2.2厘米，下2厘米，左2.7厘米，右2.3厘米。必须插入居中页码。除摘要和文献清单用单倍行距外，其余（包括正文、标题、姓名、日期）一律用1.2倍行距。中文一律用宋体，英文和数字一律用Times New Roman。摘要和文献用小五号字，正文用五号，段前段后不留空行。中文摘要段前段后各空0.5行。英文摘要段前空0.5行，段后空1行。左右适当缩进。每段文字首行缩进约2个字符。

    一级标题用小四粗体，段前空0.5行，段后不留空行。二级标题用五号粗体，左端不缩进，段前后均不留多余空格。

    报告篇幅：4000~6000字，包括图表、参考文献在内7～8页，不得超过。要学会组织材料，精简文字。要求结构合理，文字通顺，格式规范，图文并茂。以消化文献内容后自己绘制的图表为优。直接拷贝的图表要严格控制数量和大小，拒绝质量太差的图。图表应与正文文字相协调。

\section{基本要求}
    主要以口头报告内容为题，也允许适当改变。要充分阅读文献，应包括近期文献和高层次文献，要有代表性，要有足够的国际文献。要列出20篇以上，注意信息完整，格式规范统一。要在正文中标明对文献的引用。

    必须独立完成，在消化文献内容、进行分析和综合基础上，用自己的话来叙述。不应详细展开某一篇文献的细节，要根据自己的理解用简明扼要的语言进行归纳。引用必要的公式，但要控制公式数量。要有自己的见解，不得进行简单拷贝。注意：以拷贝为主的报告将被认为不合格。引用的图表和主要公式应标明文献出处。

\section{注意事项}
    围绕一个主题写文献阅读报告。可根据精读文献和泛读资料，重点讲述一个研究方向。要介绍该领域的研究意义，指出相关研究热点。可扼要介绍一种代表性方法的要点，但不能通篇只涉及一两篇文献。提出一些可研究的问题，给出自己的见解。

    \subsection{文献清单}
        列出所阅读文献的清单。给出完整信息：作者，题目，期刊，卷期，页码，年份，会议名称，会议地点，网址，学位论文出处，技术报告来源，专利号等。在报告正文中标明文献引用情况。

    \subsection{一些通病}
        文献档次低、资料陈旧、数量不足、简单应付、正文中不标注文献。以转抄拼凑网上现成资料为主，囫囵吞枣，不加取舍，缺乏见解。写成介绍某一问题或方法的介绍性文章，与文献关系甚少。仅根据一两篇资料详细讲述技术细节，照本宣科，冗长乏味。Word格式粗糙马虎，形式不规范。

    \subsection{评分依据}
        阅读文献是否充分，有无近期文献，所选文献有无代表性。是否在理解文献基础上用自己的语言进行了清晰的归纳和叙述。结构是否合理，文字是否通顺，格式是否规范。有无独到见解。口头发言水平：表达是否清楚，有条理，PPT质量，讨论情况。
        
\section{报告写作提纲}
    文献阅读书面报告应包括以下内容：
    \begin{itemize}
        \item[-] 中文题目和摘要
        \item[-] 英文题目和摘要
        \item[-] 引言：讨论什么问题，有何意义，应用领域或前景，有待解决的问题等。
        \item[-] 正文：文献综述的主要内容，可根据具体情况分节。
        \item[-] 结论：归纳报告的主要内容，可提出自己的见解。
        \item[-] 参考文献：期刊论文、会议论文、网址、学位论文、书籍、专利格式如下（注意观察）。
    \end{itemize}





















    % \input{./章节/XXX.tex}

    % \input{./章节/XXX.tex}

    % % !TEX root = ../thesis.tex

\section{结束语}
    XXX


    %- 正文结束 - -------------------------------------------------------


    %- 参考文献 - 
    % 加参考文献有两种方式 
    % 1.手动实现
    % 2.bibtex 这里默认使用bibtex 此时需要编译四次  xelatex->bibtex->xelatex*2
    % 如果没有使用bibtex 引用文献, 而后面用bibtex编译会报错
    % #region
    % \newpage % 参考文献另起一页
    \phantomsection\addcontentsline{toc}{section}{参考文献} % 把参考文献 以section级 加到目录
    \titleformat{\section}{}{\thesection}{1em}{} % 让 参考文献 几个字居中 好像没有更好的实现方法了

    % \nocite{Durrande2017} % 如果没有引用的文献也想显示在参考文献中
    \nocite{herrmann2018object}
    \nocite{herrmann2018robust}
    \nocite{li_bundle_2018}
    \nocite{zhang_parallax-tolerant_2014}
    
    \bibliography{bibcite} % 生成参考文献列表


    % 附录  
    % \newpage
    % % !TEX root = ../thesis.tex
% 附录全手动实现
\phantomsection\addcontentsline{toc}{section}{附录 A}
\titleformat{\section}{\sanhao\bfseries\centering}{\thescetion}{1em}{}
\setcounter{subsection}{0}
\setcounter{equation}{0}

\section*{附录\quad A} % 作为标题

\setcounter{section}{10} % 将section 设置为10 为了超链接正确
\renewcommand\thesection{A}

    \subsection{协方差矩阵性质} \label{A.1}

        协方差矩阵总是对称半正定的. 说明其对称是容易的, 即

        \begin{equation}
            \Sigma_{i,j}=\mathrm{cov}(Y_{i},Y_{j})=\mathrm{cov}(Y_{j},Y_{i})=\Sigma_{j,i}.
        \end{equation}

        下面证明协方差矩阵是半正定的.
        \begin{pf}
            设$X$为一服从均值向量为$\mu$,协方差矩阵为$\Sigma$的$d$维高斯随机向量,$\mu=\mathrm{E}(X)$,$\Sigma=\mathrm{var}(X)$,由协方差矩阵的定义可得
            \begin{equation}
                \Sigma=\mathrm{var}(X)=\mathrm{E}[(X-\mu)(X-\mu)^{T}].
            \end{equation}
            设$\forall\alpha\in\mathds{R}^{d}$,有
            \begin{equation}
                \alpha^{T}\Sigma\alpha=\alpha^{T}\mathrm{E}[(X-\mu)(X-\mu)^{T}]\alpha=\mathrm{E}[\alpha^{T}(X-\mu)(X-\mu)^{T}\alpha].
            \end{equation}
            令$Y=\alpha^{T}(X-\mu)=(X-\mu)^{T}\alpha$, $X-\mu$是服从于中心化后(centered)的多维高斯分布, 由多维高斯分布的定义可知$Y\sim\mathcal{N}$, 且易得$\mathrm{E}(Y)=0$,于是

            \begin{equation}
                \alpha^{T}\Sigma\alpha=\mathrm{E}(Y^{2})=\mathrm{var}(Y)\geqslant 0.
            \end{equation}
            由矩阵半正定的定义可得$\Sigma$为半正定矩阵.
        \end{pf}

    \subsection{高斯条件分布} \label{A.2}

        为证明式\ref{muc}与式\ref{sigc}, 需要先介绍分块矩阵求逆定理(inversion of a
        partitioned matrix)\cite{rasmussen2006gaussian}. 设矩阵$A$为分块矩阵, 其逆阵为$A^{-1}$

        \begin{equation}
            A=  \begin{pmatrix}
                    P & Q \\
                    R & S \\
                \end{pmatrix}, \qquad\qquad
            A^{-1}= \begin{pmatrix}
                        \tilde{P} & \tilde{Q} \\
                        \tilde{R} & \tilde{S} \\
                    \end{pmatrix}.
        \end{equation}

        其中$P$与$\tilde{P}$为$n_{1}\times n_{1}$的方阵, $S$与$\tilde{S}$为$n_{2}\times n_{2}$的方阵. 在已知$P$,$Q$,$R$,$S$的条件下, 可由下式计算得到$\tilde{P}$,$\tilde{Q}$,$\tilde{R}$,$\tilde{S}$.

        \begin{equation}
            \begin{cases}
                \tilde{P}=N. \\
                \tilde{Q}=-NQS^{-1}. \\ 
                \tilde{R}=-S^{-1}RN. \\
                \tilde{S}=S^{-1}+S^{-1}RNQS^{-1}. \label{A2.6}
            \end{cases}
        \end{equation}

        其中$N=(P-QS^{-1}R)^{-1}$. 下面在基于性质2.2的设定下完成对式\ref{muc}与式\ref{sigc}的证明.

        \begin{pf}
            不失一般性, 我们假定$Y$是经过中心化的, 即$\mu=0$. 对于高斯分布我们通常只需关注指数部分, 从多维高斯分布的pdf出发,代入\ref{2.13}我们得到
            \begin{align}\label{A2.7}
                f_{Y}(y)&\propto\mathrm{exp}\left(-\frac{1}{2}y^{T}\Sigma^{-1}y\right) \notag \\
                &\propto\mathrm{exp}\left( -\left(y^{T}_{1}\ y^{T}_{2}\right)
                \begin{pmatrix}
                    \Sigma_{1,1} & \Sigma_{1,2} \\
                    \Sigma_{2,1} & \Sigma_{2,2} \\
                \end{pmatrix}^{-1}
                \left({y_{1}\atop y_{2}}\right)\right).
            \end{align}

            下面求分块矩阵的逆,令

            \begin{equation}
                \begin{pmatrix}
                    \Sigma_{1,1} & \Sigma_{1,2} \\
                    \Sigma_{2,1} & \Sigma_{2,2} \\
                \end{pmatrix}^{-1}=
                \begin{pmatrix}
                    A & B \\
                    C & D \\
                \end{pmatrix}.
            \end{equation}

            由\ref{A2.6}计算得

            \begin{equation}\label{A2.9}
                \begin{cases}
                    A=N. \\
                    B=-N\Sigma_{1,2}\Sigma_{2,2}^{-1}. \\ 
                    C=-\Sigma_{2,2}^{-1}\Sigma_{2,1}N. \\
                    D=\Sigma_{2,2}^{-1}+\Sigma_{2,2}^{-1}\Sigma_{2,1}N\Sigma_{1,2}\Sigma_{2,2}^{-1}. 
                \end{cases} 
            \end{equation}

            上式中$N=(\Sigma_{1,1}-\Sigma_{1,2}\Sigma^{-1}_{2,2}\Sigma_{2,1})^{-1}$, 由协方差矩阵性质$\Sigma_{1,2}=\Sigma^{T}_{2,1}$.

            将式\ref{A2.9}代入式\ref{A2.7}, 并进行整理

            \begin{align}\label{A.10}
                f_{Y}(y)&\propto\mathrm{exp}\left( -\left(y^{T}_{1}\ y^{T}_{2}\right)
                \begin{pmatrix}
                    A & B \\
                    C & D \\
                \end{pmatrix}
                \left({y_{1}\atop y_{2}}\right)\right)\notag\\
                &\propto\mathrm{exp}(-y^{T}_{1}Ay_{1}-y^{T}_{2}Cy_{1}-y^{T}_{1}By_{2}-y^{T}_{2}Dy_{2} )\notag\\
                &\begin{aligned}
                    \ \propto\mathrm{exp}(  &-y^{T}_{1}(\Sigma_{1,1}-\Sigma_{1,2}\Sigma^{-1}_{2,2}\Sigma{2,1})^{-1}y_{1} \\ 
                                            &+y^{T}_{2}\Sigma^{-1}_{2,2}\Sigma_{2,1}(\Sigma_{1,1}-\Sigma_{1,2}\Sigma^{-1}_{2,2}\Sigma_{2,1})^{-1}y_{1} \\ 
                                            &+y^{T}_{1}(\Sigma_{1,1}-\Sigma_{1,2}\Sigma^{-1}_{2,2}\Sigma_{2,1})^{-1}\Sigma_{1,2}\Sigma^{-1}_{2,2}y_{2} \\ 
                                            &-y^{T}_{2}\left[\Sigma^{-1}_{2,2}+\Sigma^{-1}_{2,2}\Sigma_{2,1}(\Sigma_{1,1}-\Sigma_{1,2}\Sigma^{-1}_{2,2}\Sigma_{2,1})^{-1}\Sigma_{1,2}\Sigma^{-1}_{2,2}\right]y_{2})  \\
                \end{aligned} \notag\\
                &\begin{aligned}
                    \ \propto\mathrm{exp}(  &-y^{T}_{1}(\Sigma_{1,1}-\Sigma_{1,2}\Sigma^{-1}_{2,2}\Sigma_{2,1})^{-1}(y_{1}-\Sigma_{1,2}\Sigma^{-1}_{2,2}y_{2}) \\ 
                                            &+y^{T}_{2}\Sigma^{-1}_{2,2}\Sigma_{2,1}(\Sigma_{1,1}-\Sigma_{1,2}\Sigma^{-1}_{2,2}\Sigma_{2,1})^{-1}(y_{1}-\Sigma_{1,2}\Sigma^{-1}_{2,2}y_{2})) \\ 
                \end{aligned} \notag\\
                &\propto\mathrm{exp}(-(y^{T}_{1}-y^{T}_{2}\Sigma^{-1}_{2,2}\Sigma_{2,1})(\Sigma_{1,1}-\Sigma^{-1}_{2,2}\Sigma_{2,1})^{-1}(y_{1}-\Sigma_{1,2}\Sigma^{-1}_{2,2}y_{2})) \notag\\
                &\propto\mathrm{exp}(-(y_{1}-\Sigma_{1,2}\Sigma^{-1}_{2,2}y_{2})^{T}(\Sigma_{1,1}-\Sigma^{-1}_{2,2}\Sigma_{2,1})^{-1}(y_{1}-\Sigma_{1,2}\Sigma^{-1}_{2,2}y_{2})).
            \end{align}

            在$y_{2}$已知的条件下, $\mu_{c}=\Sigma_{1,2}\Sigma^{-1}_{2,2}y_{2}$与$\Sigma_{c}=\Sigma_{1,1}-\Sigma^{-1}_{2,2}\Sigma_{2,1}$是确定的, 故$Y_{1}|Y_{2}$服从均值为$\mu_{c}$,协方差矩阵为$\Sigma_{c}$的高斯分布, 即

            $$Y_{1}|Y_{2}\sim\mathcal{N}(\mu_{c},\Sigma_{c}).$$

            如果不考虑进行中心化, 则只需将式\ref{A.10}中$y_{1}$与$y_{2}$分别用$y_{1}-\mu_{1}$与$y_{2}-\mu_{2}$替代, 即得到式\ref{muc}与式\ref{sigc}结果.

        \end{pf}

    \subsection{贝叶斯线性模型$\boldsymbol{w}$的后验分布推导} \label{A.3}

        将式\ref{2.23}, 式\ref{2.24}代入式\ref{2.25}进行整理
        \begin{equation}\label{A3.11}
        \begin{aligned}
            p(\boldsymbol{w}|\boldsymbol{X},\boldsymbol{y})
            &\propto p(\boldsymbol{y}|\boldsymbol{X},\boldsymbol{w})p(\boldsymbol{w}) \\
            &\propto \mathrm{exp}\left(-\frac{1}{2\sigma^{2}_{n}}(\boldsymbol{y}-\boldsymbol{X}^{T}\boldsymbol{w})^{T}(\boldsymbol{y}-\boldsymbol{X}^{T}\boldsymbol{w})\right) \mathrm{exp}\left(-\frac{1}{2}\boldsymbol{w}^{T}\Sigma^{-1}_{p}\boldsymbol{w} \right) \\
            &\propto \mathrm{exp}\left(-\frac{1}{2}\left[\frac{1}{\sigma^{2}_{n}}(\boldsymbol{y}-\boldsymbol{X}^{T}\boldsymbol{w})^{T}(\boldsymbol{y}-\boldsymbol{X}^{T}\boldsymbol{w})-\boldsymbol{w}^{T}\Sigma^{-1}_{p}\boldsymbol{w}\right] \right) \\
            &\propto \mathrm{exp}\left(-\frac{1}{2}\left[ \frac{1}{\sigma^{2}_{n}}\boldsymbol{y}^{T}\boldsymbol{y} - \frac{1}{\sigma^{2}_{n}}\boldsymbol{y}^{T}\boldsymbol{X}^{T}\boldsymbol{w} - \frac{1}{\sigma^{2}_{n}}\boldsymbol{w}^{T}\boldsymbol{Xy} + \frac{1}{\sigma^{2}_{n}}\boldsymbol{w}^{T}\boldsymbol{X}\boldsymbol{X}^{T}\boldsymbol{w} + \boldsymbol{w}^{T}\Sigma^{-1}_{p}\boldsymbol{w} \right] \right) \\
            &\propto \mathrm{exp}\left(-\frac{1}{2}\left[- \frac{2}{\sigma^{2}_{n}}\boldsymbol{w}^{T}\boldsymbol{Xy} + \frac{1}{\sigma^{2}_{n}}\boldsymbol{w}^{T}\boldsymbol{X}\boldsymbol{X}^{T}\boldsymbol{w} + \boldsymbol{w}^{T}\Sigma^{-1}_{p}\boldsymbol{w} \right] \right) \\
            &\propto \mathrm{exp}\left(-\frac{1}{2}\left[ -\frac{2}{\sigma^{2}_{n}}\boldsymbol{w}^{T}\boldsymbol{Xy} + \boldsymbol{w}^{T}(\frac{1}{\sigma^{2}_{n}}\boldsymbol{X}\boldsymbol{X}^{T} + \Sigma^{-1}_{p})\boldsymbol{w} \right] \right) \\
            &\propto \mathrm{exp}\left(-\frac{1}{2}\left[ -\frac{2}{\sigma^{2}_{n}}\boldsymbol{w}^{T}\boldsymbol{Xy} + \boldsymbol{w}^{T}A\boldsymbol{w} \right] \right).
        \end{aligned}
        \end{equation}

        其中$A=\sigma^{-2}\boldsymbol{X}\boldsymbol{X}^{T} + \Sigma^{-1}_{p}$. 最后需要凑成多维高斯分布指数的形式, 可以使用待定系数法来求$\boldsymbol{\bar{w}}$. 从目标形式出发

        \begin{equation}\label{A3.12}
            \begin{aligned}
                p(\boldsymbol{w}|\boldsymbol{X},\boldsymbol{y})
                &\propto \mathrm{exp}\left( -\frac{1}{2}(\boldsymbol{w}-\boldsymbol{\bar{w}})^{T}A(\boldsymbol{w}-\boldsymbol{\bar{w}}) \right) \\
                &\propto \mathrm{exp}\left( -\frac{1}{2}(\boldsymbol{w}^{T}A\boldsymbol{w} - \boldsymbol{w}^{T}A\boldsymbol{\bar{w}} - \boldsymbol{\bar{w}}^{T}A\boldsymbol{w} + \boldsymbol{\bar{w}}^{T}A\boldsymbol{\bar{w}}) \right) \\
                &\propto \mathrm{exp}\left( -\frac{1}{2}(\boldsymbol{w}^{T}A\boldsymbol{w} - 2\boldsymbol{w}^{T}A\boldsymbol{\bar{w}}) \right). \\
            \end{aligned}
        \end{equation}

        上式推导中注意$\boldsymbol{\bar{w}}$应为与$\boldsymbol{w}$无关的常数. 对比式\ref{A3.11}与式\ref{A3.12}即可得

        \begin{equation}
            \begin{aligned}
                A\boldsymbol{\bar{w}}&=\frac{1}{\sigma^{2}_{n}}\boldsymbol{Xy} \\
                \boldsymbol{\bar{w}}&=\frac{1}{\sigma^{2}_{n}}A^{-1}\boldsymbol{Xy}.
            \end{aligned}
        \end{equation}

    \subsection{贝叶斯线性模型预测分布推导} \label{A.4}

        % 这个推导依然有疑问 这里只是我想到的一种可能性

        为了得到$f_{*}|\boldsymbol{x}_{*},\boldsymbol{X},\boldsymbol{y}$的后验分布, 先求$y_{*}|\boldsymbol{x}_{*},\boldsymbol{X},\boldsymbol{y}$的后验分布, 然后利用$y_{*}=f_{*}+\varepsilon_{*}$求得$f_{*}|\boldsymbol{x}_{*},\boldsymbol{X},\boldsymbol{y}$.

        \begin{equation} \label{A4.14}
            \begin{aligned}
            p(y_{*}|\boldsymbol{x}_{*},\boldsymbol{X},\boldsymbol{y})
            &=\int p(y_{*}|\boldsymbol{x}_{*},\boldsymbol{w})p(\boldsymbol{w}|\boldsymbol{X},\boldsymbol{y}) \mathrm{d}\boldsymbol{w}. \\
            \end{aligned}
        \end{equation}

        借助式\ref{2.22}与式\ref{2.26}, 有


        \begin{equation}
        \begin{aligned}\label{A4.15}
            p(y_{*}|\boldsymbol{x}_{*},\boldsymbol{X},\boldsymbol{y})
            &\propto\int\mathrm{exp}\left(-\frac{(y_{*}-\boldsymbol{w}^{T}\boldsymbol{x}_{*})^{2}}{2\sigma^{2}_{n}} \right)\mathrm{exp}\left( -\frac{1}{2}(\boldsymbol{w}-\boldsymbol{\bar{w}})^{T}A(\boldsymbol{w}-\boldsymbol{\bar{w}}) \right)\mathrm{d}\boldsymbol{w} \\
            &\propto\int\mathrm{exp}\left[ -\frac{1}{2}\left(\frac{1}{\sigma^{2}_{n}}y_{*}^{2} - \frac{2}{\sigma^{2}_{n}}\boldsymbol{w}^{T}\boldsymbol{x}y_{*} + \frac{1}{\sigma^{2}_{n}}\boldsymbol{w}^{T}\boldsymbol{x}_{*}\boldsymbol{x}_{*}^{T}\boldsymbol{w} + \boldsymbol{w}^{T}A\boldsymbol{w} - 2\boldsymbol{w}^{T}A\boldsymbol{\bar{w}} \right)\right] \mathrm{d}\boldsymbol{w} \\
            &\propto\int\mathrm{exp}\left[ -\frac{1}{2}\left( \boldsymbol{w}^{T}(\frac{1}{\sigma^{2}_{n}}\boldsymbol{x}_{*}\boldsymbol{x}^{T}_{*}+A )\boldsymbol{w} - 2\boldsymbol{w}^{T}(\frac{1}{\sigma^{2}_{n}}\boldsymbol{x}_{*}y_{*}+A\boldsymbol{\bar{w}} ) + \frac{1}{\sigma^{2}_{n}}y^{2}_{*} \right)\right] \mathrm{d}\boldsymbol{w} \\
            &\propto\int\mathrm{exp}\left[ -\frac{1}{2}\left( \boldsymbol{w}^{T}L\boldsymbol{w} - 2\boldsymbol{w}^{T}(\frac{1}{\sigma^{2}_{n}}\boldsymbol{x}_{*}y_{*}+A\boldsymbol{\bar{w}} ) + \frac{1}{\sigma^{2}_{n}}y^{2}_{*} \right)\right]\mathrm{d}\boldsymbol{w}. \\
        \end{aligned}
        \end{equation}

        上式中$L=\sigma^{-2}_{n}\boldsymbol{x}_{*}\boldsymbol{x}^{T}_{*}+A$, 构造$\boldsymbol{w}$服从均值为$\boldsymbol{m}$, 协方差矩阵为$L$的高斯分布指数部分, 用待定系数法求出$\boldsymbol{m}$.

        \begin{equation}\label{A4.16}
            \begin{aligned}
                &(\boldsymbol{w}-\boldsymbol{m})^{T}L(\boldsymbol{w}-\boldsymbol{m}) \\
                =&\boldsymbol{w}^{T}L\boldsymbol{w}-2\boldsymbol{w}^{T}L\boldsymbol{m}+\boldsymbol{m}^{T}L\boldsymbol{m}
            \end{aligned}
        \end{equation}

        对比式\ref{A4.15}与式\ref{A4.16}, 可求出$\boldsymbol{m}$.

        \begin{equation}\label{A4.17}
            \begin{aligned}
                L\boldsymbol{m}&=\frac{1}{\sigma^{2}_{n}}\boldsymbol{x}_{*}y_{*}+A\boldsymbol{\bar{w}} \\
                \boldsymbol{m}&=L^{-1}(\frac{1}{\sigma^{2}_{n}}\boldsymbol{x}_{*}y_{*}+A\boldsymbol{\bar{w}})
            \end{aligned}
        \end{equation}

        $\boldsymbol{m}$与$L$是独立于$\boldsymbol{w}$的, 于是

        \begin{equation}
            \begin{aligned}\label{A4.18}
                p(y_{*}|\boldsymbol{x}_{*},\boldsymbol{X},\boldsymbol{y})
                &\propto\int\mathrm{exp}\left( -\frac{1}{2}(\boldsymbol{w}-\boldsymbol{m})^{T}L(\boldsymbol{w}-\boldsymbol{m}) \right) \mathrm{exp}\left[ -\frac{1}{2}(\frac{1}{\sigma^{2}_{n}}y^{2}_{*} - \boldsymbol{m}^{T}L\boldsymbol{m}) \right]  \mathrm{d}\boldsymbol{w} \\
                &\propto\mathrm{exp}\left[ -\frac{1}{2}\left(\frac{1}{\sigma^{2}_{n}}y^{2}_{*} - \boldsymbol{m}^{T}L\boldsymbol{m}\right) \right]\int\mathrm{exp}\left( -\frac{1}{2}(\boldsymbol{w}-\boldsymbol{m})^{T}L(\boldsymbol{w}-\boldsymbol{m}) \right) \mathrm{d}\boldsymbol{w} \\
                &\propto\mathrm{exp}\left[ -\frac{1}{2}\left(\frac{1}{\sigma^{2}_{n}}y^{2}_{*} - \boldsymbol{m}^{T}L\boldsymbol{m}\right) \right]\\
            \end{aligned}
        \end{equation}

        \begin{equation}
            \begin{aligned}\label{A4.19}
                \boldsymbol{m}^{T}L\boldsymbol{m}
                &=\left(\frac{1}{\sigma^{2}_{n}}\boldsymbol{x}_{*}y_{*}+A\boldsymbol{\bar{w}}\right)^{T}L^{-1}\cdot L\cdot L^{-1}\left(\frac{1}{\sigma^{2}_{n}}\boldsymbol{x}_{*}y_{*}+A\boldsymbol{\bar{w}}\right) \\
                &=\left( \frac{1}{\sigma^{4}_{n}}\boldsymbol{x}^{T}_{*}L^{-1}\boldsymbol{x}_{*} \right)y^{2}_{*} + 2\left( \frac{1}{\sigma^{2}_{n}}\boldsymbol{x}^{T}_{*}L^{-1}A\boldsymbol{\bar{w}}\right)y_{*}+\boldsymbol{\bar{w}}^{T}AL^{-1}A\boldsymbol{\bar{w}} \\
            \end{aligned}
        \end{equation}

        把式\ref{A4.19}代入\ref{A4.18}, 注意$L$,$A$,$\boldsymbol{\bar{w}}$都独立于$y_{*}$,整理得

        \begin{equation}
            \begin{aligned} \label{A4.20}
                p(y_{*}|\boldsymbol{x}_{*},\boldsymbol{X},\boldsymbol{y})
                &\propto\mathrm{exp}\left[ -\frac{1}{2}\left(\frac{1}{\sigma^{2}_{n}}y^{2}_{*} - \boldsymbol{m}^{T}L\boldsymbol{m}\right) \right]\\
                &\propto\mathrm{exp}\left[ -\frac{1}{2}\left( \Bigl(\frac{1}{\sigma^{2}_{n}}-\frac{1}{\sigma^{4}_{n}}\boldsymbol{x}^{T}_{*}L^{-1}\boldsymbol{x}_{*} \Bigl)y^{2}_{*}-2\Bigl(\frac{1}{\sigma^{2}_{n}}\boldsymbol{x}^{T}_{*}L^{-1}A\boldsymbol{\bar{w}} \Bigl)y_{*} \right) \right] \\
            \end{aligned}
        \end{equation}

        同样从目标出发

        \begin{equation} \label{A4.21}
            \lambda(y_{*}-\bar{y}_{*})^{2}=\lambda y^{2}_{*}-2\lambda \bar{y}_{*}y_{*}+\lambda \bar{y}^{2}_{*}
        \end{equation}

        对比式\ref{A4.20}括号内部分与式\ref{A4.21}得

            \begin{align} 
                &\lambda =\frac{1}{\sigma^{2}_{n}}(1-\frac{1}{\sigma^{2}_{n}}\boldsymbol{x}^{T}_{*}L^{-1}\boldsymbol{x}_{*}). \label{A4.22}\\
                &\bar{y}_{*}=\frac{1}{\lambda}(\frac{1}{\sigma^{2}_{n}}\boldsymbol{x}^{T}_{*}L^{-1}A\boldsymbol{\bar{w}}). \label{A4.23}
            \end{align}

            于是

            \begin{equation}
                \begin{aligned} \label{A4.24}
                    p(y_{*}|\boldsymbol{x}_{*},\boldsymbol{X},\boldsymbol{y})
                    &\propto\mathrm{exp}\left( -\frac{1}{2}\lambda (y_{*} - \bar{y}_{*})^{2} \right).\\
                \end{aligned}
            \end{equation}

            下面计算$\lambda$与$\bar{y}_{*}$以得到预测分布得均值与方差. 为简便起见, 首先计算$L^{-1}$, $\boldsymbol{x}^{T}_{*}L^{-1}\boldsymbol{x}_{*}$, 以及$\boldsymbol{x}^{T}_{*}L^{-1}$. 计算$L^{-1}$需要用到矩阵求逆引理(matrix inversion lemma)\cite{rasmussen2006gaussian}. 以下计算中记$\alpha=\boldsymbol{x}^{T}_{*}A^{-1}\boldsymbol{x}$.

            \begin{equation}
                \begin{aligned} \label{A4.25}
                    L^{-1}
                    &=\left(A+\boldsymbol{x}_{*}\frac{1}{\sigma^{2}_{n}}\boldsymbol{x}^{T}_{*}\right)^{-1} \\
                    &=A^{-1}-A^{-1}\boldsymbol{x}_{*}(\sigma^{2}_{n}+\boldsymbol{x}^{T}_{*}A^{-1}\boldsymbol{x}_{*})^{-1}\boldsymbol{x}^{T}_{*}A^{-1} \\
                    &=A^{-1}-\frac{1}{\sigma^{2}_{n}+\alpha}A^{-1}\boldsymbol{x}_{*}\boldsymbol{x}^{T}_{*}A^{-1}.
                \end{aligned}
            \end{equation}

            \begin{equation}
                \begin{aligned} \label{A4.26}
                    \boldsymbol{x}^{T}_{*}L^{-1}\boldsymbol{x}_{*}
                    &=\boldsymbol{x}^{T}_{*}A^{-1}\boldsymbol{x}_{*}-\frac{1}{\sigma^{2}_{n}+\alpha}\boldsymbol{x}^{T}_{*}A^{-1}\boldsymbol{x}_{*}\boldsymbol{x}^{T}_{*}A^{-1}\boldsymbol{x}_{*}=\frac{\alpha \sigma^{2}_{n}}{\sigma^{2}_{n}+\alpha}.
                \end{aligned}
            \end{equation}

            \begin{equation}\label{A4.27}
                    \boldsymbol{x}^{T}_{*}L^{-1}=\frac{\sigma^{2}_{n}}{\sigma^{2}_{n}+\alpha} \boldsymbol{x}^{T}_{*}A^{-1}.
            \end{equation}

            故

            \begin{equation}\label{A4.28}
                \lambda=\frac{1}{\sigma^{2}_{n}+\alpha}.
            \end{equation}

            \begin{equation}\label{A4.29}
                \bar{y}_{*}=\frac{1}{\sigma^{2}_{n}}\boldsymbol{x}^{T}_{n}A^{-1}\boldsymbol{Xy}.
            \end{equation}

            \begin{equation}\label{A4.30}
                y_{*}|\boldsymbol{x}_{*},\boldsymbol{X},\boldsymbol{y}\sim\mathcal{N}(\frac{1}{\sigma^{2}_{n}}\boldsymbol{x}^{T}_{n}A^{-1}\boldsymbol{Xy}, \frac{1}{\lambda}=\sigma^{2}_{n}+\boldsymbol{x}^{T}_{*}A^{-1}\boldsymbol{x}).
            \end{equation}

            由$f_{*}=y_{*}-\varepsilon_{*}$

            \begin{equation}\label{A4.31}
                f_{*}|\boldsymbol{x}_{*},\boldsymbol{X},\boldsymbol{y}\sim\mathcal{N}(\frac{1}{\sigma^{2}_{n}}\boldsymbol{x}^{T}_{n}A^{-1}\boldsymbol{Xy}, \boldsymbol{x}^{T}_{*}A^{-1}\boldsymbol{x}).
            \end{equation}
            
    \subsection{贝叶斯线性模型预测分布变形} \label{A.5}
        首先对式\ref{2.30}均值进行变形.

        \begin{equation} \label{A5.32}
            \begin{aligned}
                \boldsymbol{\phi}^{T}_{*}\left( \frac{1}{\sigma^{2}_{n}}A^{-1}\boldsymbol{\varPhi} \right)\boldsymbol{y}
                &=\boldsymbol{\phi}^{T}_{*}\left[ A^{-1}\cdot \frac{1}{\sigma^{2}_{n}}\boldsymbol{\varPhi}(K+\sigma^{2}_{n}I)\cdot (K+\sigma^{2}_{n}I)^{-1} \right] \boldsymbol{y} \\
                &=\boldsymbol{\phi}^{T}_{*}\left[ A^{-1}\cdot A\Sigma_{p}\boldsymbol{\varPhi} \cdot (K+\sigma^{2}_{n}I)^{-1} \right] \boldsymbol{y} \\
                &=\boldsymbol{\phi}^{T}_{*}\Sigma_{p}\boldsymbol{\varPhi}(K+\sigma^{2}_{n}I)^{-1}\boldsymbol{y}.
            \end{aligned}
        \end{equation}

        然后方差的变形需要先用矩阵求逆引理展开$A^{-1}$.

        \begin{equation} \label{A5.33}
            \begin{aligned}
                A^{-1}&=\left(\frac{1}{\sigma^{2}_{n}}\boldsymbol{\varPhi}\boldsymbol{\varPhi}^{T}+\Sigma^{-1}_{p}\right)^{-1} \\
                &=\Sigma_{p}-\Sigma_{p}\boldsymbol{\varPhi}(\boldsymbol{\varPhi}^{T}\Sigma_{p}\boldsymbol{\varPhi} + \sigma^{2}_{n}I)^{-1}\boldsymbol{\varPhi}^{T}\Sigma_{p} \\
                &=\Sigma_{p}-\Sigma_{p}\boldsymbol{\varPhi}(K + \sigma^{2}_{n}I)^{-1}\boldsymbol{\varPhi}^{T}\Sigma_{p} \\
            \end{aligned}
        \end{equation}

        故

        \begin{equation} \label{A5.34}
                \boldsymbol{\phi}_{*}^{T}A^{-1}\boldsymbol{\phi}_{*}=\boldsymbol{\phi}^{T}_{*}\Sigma_{p}\boldsymbol{\phi}_{*} -  \boldsymbol{\phi}^{T}_{*}\Sigma_{p}\boldsymbol{\varPhi}(K+\sigma^{2}_{n}I)^{-1} \boldsymbol{\varPhi}^{T}\Sigma_{p}\boldsymbol{\phi}_{*}
        \end{equation}

    
    % %- 参考文献 - 不使用bibtex的方式  手动实现
    % \phantomsection\addcontentsline{toc}{part}{参考文献} % 把参考文献 以part级 加到目录 如果使用bibtex来处理文献不用这个
    % %\titleformat{\section}{\centering}{\thesection}{1em}{}%让 参考文献 几个字居中 好像没有更好的实现方法了
    % \begin{thebibliography}{99}% “99” 表示以最多两位数来编号参考文献，用于对齐
    %     \addtolength{\itemsep}{-2ex} % 用于更改文献行距
    %     % 文献列表
    %     \bibitem{1}author.title[EB/OL].

    % \end{thebibliography}

    % #endregion
    

\end{document}
% end  文档完成  *^_^*